% Template for post or presentation abstracts submitted to ILAS 2022
% 1. Fill in the presenter's information (name, affiliation, email
% address) and talk information (title of presentation, abstract).
% 2. Compile to make sure there are no errors.
% 3. Save the .tex file for your abstract with a distinctive name,
% ideally "FamilynameGivenname.tex".
% 4. Upload your .tex file to https://clr.ie/129557
%
% * Please keep your LaTeX commands simple, and avoid the use of
%    user-defined macros.
% * Include details of co-authors and funding acknowledgements after the abstract.
% * Upload only the LaTeX file (with .tex extension).
% * Questions: Contact Niall Madden (Niall.Madden@NUIGalway.ie)

\documentclass[a4paper]{amsart}
\usepackage{amssymb}
\usepackage{hyperref}
\pagestyle{empty}
\date{}

%% IGNORE THE NEXT 4 LINES
\makeatletter %
\def\labelenumi{\theenumi.}
\def\theenumi{\@arabic\c@enumi}
\makeatother
%% STOP IGNORING NOW

\begin{document}

\title{Barycenters of Hellinger distances and Kubo-Ando means as barycenters}
\author{D\'aniel Virosztek} %Presenting author only
\address{Alf\'ed R\'enyi Institute of Mathematics, Hungary} % Affiliation of presenting author
\email{virosztek.daniel@renyi.hu} % Email address of presenting author only

\maketitle

The first part of the talk is devoted to quantum Hellinger distances --- introduced recently by Bhatia et al. \cite{bhatia-paper} --- with a particular emphasis on barycenters. We introduce the family of generalized quantum Hellinger divergences that are of the form $\phi(A,B)=\mathrm{Tr} \left((1-c)A + c B - A \sigma B \right),$ where $\sigma$ is an arbitrary Kubo-Ando mean, and $c \in (0,1)$ is the weight of $\sigma.$ We note that these divergences belong to the family of maximal quantum $f$-divergences, and hence are jointly convex, and satisfy the data processing inequality (DPI). We will present a fixed-point equation that characterizes of the barycenter of finitely many positive definite operators for these generalized quantum Hellinger divergences \cite{pitrik-virosztek-hellinger}.
\par
In the second part, we present a divergence center interpretation of general symmetric Kubo-Ando means \cite{pitrik-virosztek-KA-bary}.
This characterization of the symmetric means naturally leads to a definition of weighted and multivariate versions of a large class of symmetric Kubo-Ando means. We study elementary properties of these weighted multivariate means, and note in particular that in the special case of the geometric mean we recover the weighted $\mathcal{A} \# \mathcal{H}$-mean introduced by Kim, Lawson, and Lim \cite{kim-lawson-lim}.

% The abstract  ***no more than one page***


%% Coauthor details here (optional - delete if not applicable)
% and funding acknowledgment (optional - delete if not applicable)
\emph{This is joint work with J\'ozsef Pitrik (TU Budapest).
Virosztek is supported by the Momentum program of the Hungarian Academy of Sciences under grant agreement no. LP2021-15/2021.}

\begin{thebibliography}{1}
\bibitem{bhatia-paper}
R. Bhatia, S. Gaubert, and T. Jain.
\newblock Matrix versions of the Hellinger distance.
\newblock {\em  Lett. Math. Phys.}, 109:1777-1804, 2019.

\bibitem{kim-lawson-lim}
S. Kim, J. Lawson, and Y. Lim. 
\newblock The matrix geometric mean of parametrized, weighted arithmetic and harmonic means.
\newblock {\em Linear Algebra Appl.}, 435:2114--2131, 2011.

\bibitem{pitrik-virosztek-hellinger}
J. Pitrik and D. Virosztek.
\newblock Quantum Hellinger distances revisited.
\newblock {\em Lett. Math. Phys.}, 110:2039-2052, 2020.

\bibitem{pitrik-virosztek-KA-bary}
J. Pitrik and D. Virosztek.
\newblock A divergence center interpretation of general symmetric Kubo-Ando means, and related weighted multivariate operator means.
\newblock {\em Linear Algebra Appl.}, 609:203-217, 2021.

\end{thebibliography}


\end{document}


% Please leave the next bit alone
%%% Local Variables:
%%% mode: latex
%%% TeX-master: "../ILAS-2022"
%%% End:

